\documentclass[12pt,a4paper]{article}

\usepackage[T2A]{fontenc}
\usepackage[utf8]{inputenc}
\usepackage[russian,english]{babel}
\usepackage{amsmath,amssymb}
\usepackage{geometry}
\usepackage{indentfirst}

\geometry{margin=2.5cm}

\begin{document}

\begin{center}
    {\LARGE Текст к конференции}
\end{center}

\section*{Содержание (7 мин)}
\begin{enumerate}
    \item Введение (2 мин)
    \item Программная имплементация (1 мин)
    \item Численный анализ модели Дикманна--Лоу (2 мин)
    \item Сравнение с \textit{neighborhood-dependent} моделью (2 мин)
\end{enumerate}

\section*{Введение (2 мин)}
Модель популяционной динамики в биологии --- это математическое описание, используемое для прогнозирования изменений численности, структуры и плотности популяции во времени. Предположим, у нас есть площадь-поле $N \times N$, заполненное индивидами, и дискретный временной шаг $\Delta t$. Каждый индивид размножается, то есть создаёт свою копию на поле, и умирает, то есть исчезает из поля, с определённой вероятностью.

Существуют несколько подходов к вычислению популяционной динамики при заданных параметрах; в данной работе мы рассматриваем модель Дикманна--Лоу, описанную в 2000 году, и \textit{neighborhood-dependent} модель Эрнандес-Гарсия и Лопеса, описанную в 2005 году.

В \textit{neighborhood-dependent} модели (далее --- ND) параметр рождаемости зависит от количества соседей в окружении индивида при заданном радиусе. Полученная сумма индивидов делится на параметр чувствительности к плотности и вычитается из базового параметра рождаемости:
\[
\lambda_i = \lambda_0 - \frac{1}{N_s} N_R^{(i)}.
\]
Смертность задаётся константным значением. В ND-модели индивиды двигаются по координатам $x$ и $y$ в соответствии с броуновским движением с параметром диффузии $D$.

В модели Дикманна--Лоу, наоборот, параметр рождаемости задаётся константой, а смертность для каждого индивида высчитывается через так называемое гладкое ядро взаимодействия:
\[
d_i = d_0 + \alpha C_i,
\qquad
C_i = \sum_j \exp\left(-\frac{r_{ij}^2}{2R_c^2}\right),
\]
где $R_c$ --- характеристический радиус взаимодействия. Полученное ядро взаимодействия умножается на параметр чувствительности к конкуренции $\alpha$ и складывается с базовой смертностью $d_0$.

Заметим, что в отличие от ND-модели модель Дикманна--Лоу не предусматривает броуновского движения особей. Тем не менее потомки во время рождения перемещаются от родителя на расстояние, которое задаётся нормальным шумом, например
\[
x_{\text{child}} = x_i + \xi,
\qquad
\xi \sim \mathcal{N}(0, \sigma^2),
\]
и аналогично по второй координате. В ND-модели потомки, напротив, появляются в том же месте, где находился родитель.

\section*{Программная имплементация (1 мин)}
Данные пространственные стохастические модели были имплементированы через объектно-ориентированное программирование на языке C++, визуализация в свою очередь создана через скрипты на Python, которые используют библиотеку \texttt{plotly}.

Каждый индивид представлен через структуру, у которой есть идентификатор, координаты и статус (жив или мёртв). Случайные числа в проекте генерируются через вихрь Мерсенна. После запуска симуляции результаты записываются в \texttt{csv}-таблицу, где в каждой строке фиксируются время, идентификатор индивида, координаты, статус и идентификатор родителя.

Асимптотики моделей ND и Дикманна--Лоу совпадают: обе обладают верхним порогом количества процессов порядка $N^2$, где $N$ --- количество особей. Реальная скорость симуляции очень сильно зависит от самой динамики популяции; в худшем случае, например когда рождаемость очень высокая, речь идёт об экспоненциальной сложности.

\section*{Численный анализ модели Дикманна--Лоу (2 мин)}
Были проведены три серии симуляций модели Дикманна--Лоу: зависимость динамики от характеристического радиуса взаимодействия и разброса потомков, зависимость динамики от параметра чувствительности к конкуренции $\alpha$, а также зависимость динамики от базовой рождаемости и смертности.

В первой серии были выбраны маленький и большой характеристический радиус взаимодействия вместе с маленьким и большим разбросом потомков. Было обнаружено, что при меньшем разбросе и при большем радиусе взаимодействия индивиды склонны собираться в кластеры, в то время как в обратном случае формируется однородное поле.

Во второй серии рассматривалось поведение индивидов при разном коэффициенте ядра $\alpha$ и других фиксированных параметрах, когда модель была настроена на кластеризацию. При маленьком коэффициенте конкуренции, стремящемся к нулю, рост популяции является экспоненциальным. При высоком коэффициенте конкуренции модель падает до определённого низкого значения, но не вымирает.

При этом в ходе исследований было выявлено стационарное значение --- это значение, при котором кумулятивное количество родившихся индивидов равно количеству умерших индивидов за всё время, причём рост обоих значений линеен. Это происходит при $\alpha \approx 0.00258$. Все значения выше данного приводят к динамике, схожей с режимом сильной конкуренции, а все значения ниже --- к динамике слабой конкуренции.

Наконец, в третьей серии мы рассматриваем поведение индивидов при разных значениях рождаемости и смертности. Заметим, что даже когда рождаемость превышает смертность, популяция всё ещё стремится к вымиранию --- сказывается влияние кластеризации. Было выявлено, что критическая разница при заданных параметрах равна $0.042$. При этом поведение при рождаемости $0.40$ и смертности $0.38$ оказывается более радикальным и сильнее склонным к вымиранию, чем поведение при $0.22$ и $0.20$.

\section*{Сравнение с \textit{neighborhood-dependent} моделью (2 мин)}
Здесь были проведены четыре симуляции пары моделей \textit{neighborhood-dependent} и Дикманна--Лоу.

В ходе первой пары симуляций модели Дикманна--Лоу и ND были настроены так, чтобы они выдавали ``слабую структуру'' популяции. Было обнаружено, что за счёт индикаторной функции в определении ND-модели кластеры намного более устойчивы, чем при гладкой оболочке ядра. Примечательна также динамика живых индивидов: в ND-модели количество живых индивидов быстро упирается в потолок, в то время как в модели Дикманна--Лоу, хотя и наблюдается определённое ограничивающее значение, рост популяции происходит более гладко.

Во второй паре симуляций модели были настроены на кластеризацию. Здесь заметны чёткость кластеров при ND и больший разброс кластеров в модели Дикманна--Лоу. Примечателен тот факт, что при взгляде на количество живых индивидов во времени у Дикманна--Лоу можно заметить, что стабильный рост не гарантируется. Связано это с тем, что в данной модели количество индивидов повышает смертность, в то время как в ND оно лишь замедляет рождаемость.

В третьей паре симуляций тестируются механизмы замедления роста популяции. Так, модель Дикманна--Лоу настраивается на разреженность через больший разброс потомков и меньший радиус взаимодействия, и видно, что при достижении критического значения динамика начинает вести себя нестабильно, но в среднем не превышает определённый уровень. Теперь настроим ND-модель на больший радиус взаимодействия и заметим, что мы получаем те же кластеры, но более отдалённые друг от друга. По динамике видно, что количество индивидов быстро достигает некоторого значения и далее флуктуирует вокруг него.

Наконец, в четвёртой паре рассматривается поведение моделей в случае границы выживания снизу. Оба графика динамики медленно стремятся к нулю, однако можно заметить нестабильный характер ND в сравнении с более гладким характером кривой для модели Дикманна--Лоу.

\end{document}